\documentclass{article}
\usepackage{amsmath, amssymb, amsthm}
\usepackage{hyperref}
\usepackage{geometry}
\geometry{margin=1in}

\title{Erd\H{o}s Problem \#566}
\date{Last updated: 2025-08-31}
\author{Source: erdosproblems.com}

\begin{document}
\maketitle

\noindent\textbf{Status:} OPEN \quad \textbf{No prize}

\noindent\textbf{Tags:} graph theory, ramsey theory

\medskip
\noindent\textbf{Problem Statement:}

Let $G$ be such that any subgraph on $k$ vertices has at most $2k-3$ edges. Is it true that, if $H$ has $m$ edges and no isolated vertices, then\[R(G,H)\ll m?\]

\bigskip
\noindent\textbf{Remarks:}

In other words, is $G$ Ramsey size linear? This fails for a graph $G$ with $n$ vertices and $2n-2$ edges (for example with $H=K_n$). Erd\H{o}s, Faudree, Rousseau, and Schelp \cite{EFRS93} have shown that any graph $G$ with $n$ vertices and at most $n+1$ edges is Ramsey size linear.

Implies [567].

This problem is #31 in Ramsey Theory in the graphs problem collection.

\bigskip
\noindent\textbf{References:}

[EFRS93] Erd\H{o}s, Paul and Faudree, R. J. and Rousseau, C. C. and Schelp, R. H., Ramsey size linear graphs. Combin. Probab. Comput. (1993), 389-399.

\bigskip
\noindent\small{Source: \url{https://www.erdosproblems.com/566}}
\end{document}
